%% ============================================================
%% Academic Journal Paper Template
%% Compatible with: IEEE, ACM, Elsevier, Springer conventions
%% ============================================================

\documentclass[10pt, twocolumn]{article}

%% ---- Core Packages ----------------------------------------
\usepackage[T1]{fontenc}
\usepackage[utf8]{inputenc}
\usepackage{lmodern}
\usepackage{lettrine}
\usepackage{microtype}           % Microtypographic refinements

%% ---- Page Geometry ----------------------------------------
\usepackage[
  top=1in, bottom=1in,
  left=0.75in, right=0.75in,
  columnsep=0.25in
]{geometry}

%% ---- Mathematics ------------------------------------------
\usepackage{amsmath, amssymb, amsthm}
\usepackage{mathtools}
\usepackage{bm}                  % Bold math symbols

%% ---- Figures & Tables -------------------------------------
\usepackage{graphicx}
\usepackage{booktabs}            % Professional table rules
\usepackage{multirow}
\usepackage{array}
\usepackage{caption}
\usepackage{subcaption}
\usepackage{float}

%% ---- Algorithms -------------------------------------------
\usepackage[ruled, vlined, linesnumbered]{algorithm2e}

%% ---- Code Listings ----------------------------------------
\usepackage{listings}
\usepackage{xcolor}

\lstdefinestyle{codestyle}{
  backgroundcolor=\color{gray!8},
  basicstyle=\ttfamily\footnotesize,
  breakatwhitespace=false,
  breaklines=true,
  captionpos=b,
  commentstyle=\color{green!50!black},
  keywordstyle=\color{blue}\bfseries,
  stringstyle=\color{orange!70!black},
  numberstyle=\tiny\color{gray},
  numbers=left,
  numbersep=5pt,
  showstringspaces=false,
  frame=single,
  rulecolor=\color{gray!40},
  tabsize=2
}
\lstset{style=codestyle}

%% ---- Hyperlinks -------------------------------------------
\usepackage[
  colorlinks=true,
  linkcolor=blue!70!black,
  citecolor=green!50!black,
  urlcolor=blue!60!black
]{hyperref}
\usepackage{url}

%% ---- Bibliography -----------------------------------------
\usepackage[numbers, sort&compress]{natbib}

%% ---- Miscellaneous ----------------------------------------
\usepackage{lipsum}              % Lorem ipsum filler — remove in production
\usepackage{enumitem}
\usepackage{siunitx}             % SI units: \SI{3.14}{\mega\hertz}
\usepackage{cleveref}            % Smart cross-references: \cref{fig:foo}

%% ---- Theorem Environments ---------------------------------
\theoremstyle{plain}
\newtheorem{theorem}{Theorem}[section]
\newtheorem{lemma}[theorem]{Lemma}
\newtheorem{corollary}[theorem]{Corollary}
\newtheorem{proposition}[theorem]{Proposition}

\theoremstyle{definition}
\newtheorem{definition}[theorem]{Definition}
\newtheorem{example}[theorem]{Example}

\theoremstyle{remark}
\newtheorem{remark}[theorem]{Remark}

%% ---- Custom Commands --------------------------------------
\newcommand{\R}{\mathbb{R}}
\newcommand{\N}{\mathbb{N}}
\newcommand{\Z}{\mathbb{Z}}
\newcommand{\E}{\mathbb{E}}
\newcommand{\Prob}{\mathbb{P}}
\newcommand{\norm}[1]{\left\lVert #1 \right\rVert}
\newcommand{\abs}[1]{\left\lvert #1 \right\rvert}
\newcommand{\inner}[2]{\left\langle #1,\, #2 \right\rangle}
\newcommand{\ie}{\textit{i.e.}\xspace}
\newcommand{\eg}{\textit{e.g.}\xspace}
\newcommand{\etal}{\textit{et al.}\xspace}

%% ---- Caption Formatting -----------------------------------
\captionsetup{
  font=small,
  labelfont=bf,
  format=hang,
  justification=justified
}

%% ---- Header / Footer --------------------------------------
\usepackage{fancyhdr}
\pagestyle{fancy}
\fancyhf{}
\fancyhead[L]{\small\itshape Ne.app, Vol.~2, No.~1, March 2026}
\fancyhead[R]{\small\thepage}
\renewcommand{\headrulewidth}{0.4pt}

%% ============================================================
%%  FRONT MATTER
%%  ============================================================

\title{%
  \vspace{-1.5em}%
  \large\textbf{A CONCEPT OF OPERATION CODES FOR FINITE STATE MACHINES.}%
}

\author{%
  \textbf{Amlal El Mahrouss}$^{1}$\thanks{Corresponding author. Email: amlal@ne-app.eu}
  \\amlal@ne-app.eu
}

\date{%
  \small 2 April 2026, Ne.app
}

%% ============================================================
\begin{document}
%% ============================================================

\twocolumn[{%
  \maketitle
  \vspace{-0.5em}
  \rule{\linewidth}{0.4pt}
  \vspace{1em}
}]

\section{INTRODUCTION}
\lettrine[lines=1,loversize=-.4]{T}{HE} recent development of software development makes one push back and think, how could one party audit and report on vulnerable part of a finite state machine operation codes?

\subsection{THE MEASURE FUNCTION}

To measure a function, the following $\operatorname{Measure}(program\_counter, limit)$ describes a measurement of a finite state machine program from $program\_counter$ to its $limit$ variable. Assuming a linear execution pathway, we get the following theorem that for such Measurement, each program from which we know, program counter, operation codes, and limit, there is a possibility of audit-ability. Let's define `audit-ability' as the ability to measure operation codes from program counter to limit.

\section{ENCODING OPERATION CODES}

\lettrine[lines=1,loversize=-.4]{E}{NCODING} operation codes would be consisting taking from a set $S_{opcode}$ in index $n$ to execute $t$ operation a finite state machine $\operatorname{fsm}(S_{opcode}, start)$. 

\subsection{AN ENCODING TO REPRESENT A CURRENT MACHINE STATE.}

\lettrine[lines=1,loversize=-.4]{C}{ONSIDER} the following 

\begin{equation}
    \operatorname{Box}(fsm) \coloneqq \int_{start}^{end} fsm(t) \space dt
\end{equation}

Define a box as a set of encoded instructions written in a linear purpose, thus $fsm$ shall be linear and defined in between $start$ and $end$.

\section{COMMINATORY STATE MACHINES}

Now the Box could be combined to guarantee time-traveled computation, in a fashion that the program counter $t$ is never ahead of a master state machine, which in that case is the lead finite state machine of a set of state machines of length $n$.

\subsection{AN EXAMPLE OF MACHINE}

\lettrine[lines=1,loversize=-.4]{P}{ICTURE} a set of state machines as dimensions of a computable set of numbers, where any operations defined by traits in that set is authorized, however unauthorized operations would be rejected by each FSM and halted by the MSM.

\section{CONCLUSION}

\lettrine[lines=1,loversize=-.4]{C}{ONCLUDING} the paper would be by generalizing FSM with concepts, orchestration, and traits, we could theorize further on concepts where such needs would be appreciated for a better understanding of a concept or problem. 

\end{document}